\documentclass[11pt]{article}

\usepackage[margin=1in]{geometry}
\usepackage{amsmath,amssymb,amsthm}
\usepackage{mathtools}
\usepackage{algorithm}
\usepackage{algpseudocode}
\usepackage{hyperref}
\usepackage{booktabs}
\usepackage{xcolor}
\usepackage{enumitem}

\newtheorem{theorem}{Theorem}
\newtheorem{lemma}{Lemma}
\newtheorem{proposition}{Proposition}
\newtheorem{definition}{Definition}
\newtheorem{corollary}{Corollary}
\newtheorem{remark}{Remark}

\newcommand{\Aut}{\mathcal{A}}
\newcommand{\Q}{Q}
\newcommand{\Phigroup}{\Phi}
\newcommand{\Stab}{\mathrm{Stab}}
\newcommand{\CPre}{\mathrm{CPre}}
\newcommand{\PreHat}{\widehat{\mathrm{Pre}}}
\newcommand{\down}[1]{{\downarrow}#1}
\newcommand{\dom}{\succeq}
\newcommand{\type}{\mathrm{type}}

\title{Symmetry Reduction for Antichain-Based Reactive Synthesis:\\
Detection, Canonical Representation, and a Sound\\
Under-Approximating Controllable-Predecessor Algorithm}
\author{Design document --- \texttt{acacia-bonsai}, branch \texttt{optimize-vs-ltlsynt}}
\date{}

\begin{document}
\maketitle

\begin{abstract}
Acacia-Bonsai decides LTL realizability by computing, for a universal
co-B\"uchi automaton $\Aut$ obtained from the specification, the greatest
fixed point of a controllable-predecessor (\CPre) operator over
downward-closed sets (\emph{downsets}) of bounded integer vectors indexed by
the states of $\Aut$. For specifications with structural symmetry under
permutation of interchangeable components (e.g.\ $n$ arbitration clients),
the antichain representing this downset grows exponentially in $n$ even
though the automaton itself grows only linearly, and we show empirically
that this blow-up persists under every representation we tried, including a
binary-decision-diagram encoding. We derive and validate, layer by layer, a
symmetry-reduction scheme that (i) automatically \emph{detects and
structurally verifies} the permutation group under which a given
specification is invariant, directly on the automaton rather than the
formula, so that no unsound symmetry is ever assumed; (ii) represents the
downset canonically as \emph{count-vectors} over client-value \emph{types},
collapsing an exponential antichain to one of size linear in $n$ for the
symmetric arbiter family; (iii) implements the three downset operations
(\emph{contains}, \emph{union}, \emph{intersect}) on this representation,
two of them exactly and the third as a deliberately incomplete but
\emph{structurally sound} under-approximation; and (iv) assembles these into
a \CPre\ step that never enumerates a client permutation group nor a
raw antichain of exponential size. We give the soundness argument in full:
because Acacia's fixed-point computation is monotonically shrinking (a
greatest fixed point), every approximation introduced is one-directional
(the computed downset is always a subset of the true one), so every
\textsc{Realizable}/\textsc{Unrealizable} verdict the algorithm reaches
remains correct; the only cost of incompleteness is that the algorithm may
answer \textsc{Unknown} more often than an exact computation would. All
components below have been implemented and unit-tested against brute-force
ground truth or against the real (non-deterministic) automata that motivate
this work; the decision procedure is also wired into the solver as an
opt-in, bounded quotient path with unconditional fallback to the existing
antichain solver. The remaining work is performance engineering and
strategy extraction, scoped in \S\ref{sec:status}.
\end{abstract}

\tableofcontents

\section{Motivation and Empirical Setting}
\label{sec:motivation}

\subsection{The problem}
Acacia-Bonsai reduces LTL realizability to a $K$-bounded safety game solved
by a backward antichain fixed point (the \emph{Acacia} algorithm family; see
the tool paper for the base algorithm). For a family of specifications
parameterized by a number of interchangeable clients $n$ (round-robin and
priority arbiters, request/grant protocols, \dots), we observed the
following on the SYNTCOMP arbiter benchmarks (automaton state count is
$4n+3$, i.e.\ \emph{linear} in $n$; timings for \texttt{arbiterN} on the
current best acacia configuration versus \texttt{ltlsynt}):

\begin{center}
\begin{tabular}{@{}lrrrrrr@{}}
\toprule
$n$ & 5 & 6 & 7 & 8 & 9 & 10\\
\midrule
antichain size (converged) & 743 & 2439 & --- & --- & --- & ---\\
acacia wall time & fast & slow & timeout & timeout & timeout & timeout\\
\texttt{ltlsynt} wall time & 0.03s & 0.10s & 0.59s & 3.78s & 24.5s & timeout\\
\bottomrule
\end{tabular}
\end{center}

Both tools are exponential in $n$ on this family; \texttt{ltlsynt} simply
has a much better constant, buying it roughly three additional client-sizes
before it too times out. Since the automaton itself is linear-sized, the
combinatorial blow-up lives entirely in the $K$-bounded game's
\emph{winning-region antichain}, not in the automaton.

\subsection{A negative result: binary decision diagrams do not help}
\label{sec:bdd-negative}
The natural first hypothesis is that the antichain's blow-up is
\emph{representational} --- many near-duplicate, permutation-related
vectors that a shared, canonical structure such as a BDD could compress ---
rather than information-theoretically large. We tested this directly: we
dumped the converged antichains for \texttt{arbiter3}--\texttt{arbiter6} and
built a BDD (BuDDy) of their downward closure, bit-blasting each
coordinate's counter value and using dynamic variable reordering (sifting)
to find a near-optimal variable order. The result is negative and clean:

\begin{center}
\begin{tabular}{@{}rrrrr@{}}
\toprule
$n$ & antichain vectors & sifted BDD nodes & nodes/vectors\\
\midrule
3 & 45   & 211  & 4.69\\
4 & 189  & 645  & 3.41\\
5 & 743  & 2111 & 2.84\\
6 & 2439 & 6257 & 2.57\\
\bottomrule
\end{tabular}
\end{center}

Both the antichain size and the BDD node count grow exponentially at
\emph{the same base} ($\approx 3\times$ per client), with the BDD
consistently using \emph{more} nodes than there are antichain elements. No
variable ordering linearizes the structure: the ``middle section'' of each
vector genuinely carries $\Theta(n)$ bits of information about which client
holds which of a bounded number of counter values, subject to mutual
exclusion and fairness constraints, and this information is not compressible
by node sharing. \textbf{A BDD-backed downset would push the memory wall out
by only a constant factor, not break the exponential.}

\subsection{A positive result: the blow-up is exactly a symmetric-group orbit-counting problem}
\label{sec:orbit-positive}
Arbitration clients are \emph{interchangeable}: the specification (and
hence the winning region) is invariant under any permutation of client
indices. We tested whether \emph{canonicalizing} the antichain under this
symmetry --- keeping one representative per orbit of the permutation action
--- collapses the blow-up, again by direct measurement on the same dumped
antichains. Column-signature analysis first confirmed the expected block
structure: each vector's coordinates split into exactly $4$ blocks of size
$n$ (one automaton state per client, per structural role) plus $3$
client-independent (``shared'') coordinates, i.e.\ $4n+3$ total, matching
the automaton's state count exactly. We then verified, empirically, that
\emph{every} pairwise client transposition preserves the antichain (i.e.\
the full symmetric group $S_n$ is a genuine automorphism group of the
winning region) and computed orbit counts:

\begin{center}
\begin{tabular}{@{}rrrrrr@{}}
\toprule
$n$ & antichain vectors & valid client transpositions & orbits & collapse factor\\
\midrule
3 & 45   & $3/3$ (full $S_3$)  & 13 & 3.5$\times$\\
4 & 189  & $6/6$ (full $S_4$)  & 20 & 9.4$\times$\\
5 & 743  & $10/10$ (full $S_5$) & 27 & 27.5$\times$\\
6 & 2439 & $15/15$ (full $S_6$) & 34 & 71.7$\times$\\
\bottomrule
\end{tabular}
\end{center}

The orbit count is \emph{linear} in $n$ ($7n-8$, exactly), while the raw
antichain is exponential ($\approx 3^n$). This is the result the BDD
experiment could not deliver: a genuinely sub-exponential --- here, linear
--- representation of the winning region, obtained purely by exploiting the
permutation symmetry rather than by attempting node sharing on an
unstructured encoding. The remainder of this document derives a sound
algorithm that computes with this canonical representation directly,
without ever materializing the exponential raw antichain or enumerating a
permutation group of size $n!$.

\section{Background: the $K$-Bounded Safety Game}
\label{sec:background}

Fix a universal co-B\"uchi automaton $\Aut = (\Q, q_0, \delta, \mathrm{Acc})$
over an alphabet partitioned into input letters $I$ and output letters $O$,
obtained from the (possibly negated) LTL specification. A \emph{vector} is a
function $v : \Q \to \{-1, 0, \dots, K\}$; $-1$ denotes ``no obligation
remaining'' and each non-$(-1)$ value is a bound on how many more times an
accepting state may be visited before the run must be declared losing for
the safety abstraction. Vectors are ordered pointwise
($u \dom v \iff \forall q.\, u(q) \geq v(q)$), and a \emph{downset} is a set
closed downward under $\dom$; every downset is determined by its antichain
of $\dom$-maximal elements.

For a fixed input $i \in I$ and output $o \in O$ compatible with $i$, the
\emph{backward action} $\PreHat(\cdot, i, o) : \{-1,\dots,K\}^\Q \to
\{-1,\dots,K\}^\Q$ propagates a vector one step backward along the
transitions labelled $(i,o)$, decrementing on visits to accepting states and
flooring at $-1$ (this is the operator implemented, unchanged, by
\texttt{actioners::standard::apply} in the codebase; we reuse it verbatim
throughout this document and do not re-derive or modify it). The
controllable predecessor of a downset $F$ is
\[
  \CPre(F) \;=\; F \,\cap\, \bigcap_{i \in I} \bigcup_{o\ \mathrm{compat.}\ i} \PreHat(F, i, o),
\]
and Acacia computes the greatest fixed point of $\CPre$ starting from the
$K$-bounded ``safe'' downset, incrementing $K$ when the initial vector
$\iota$ (with $\iota(q_0)=0$, $\iota(q)=-1$ otherwise) drops out of the
current iterate. The specification is realizable iff this process converges
with $\iota$ inside the fixed point for some $K \leq K_{\max}$. Because the
fixed point is computed as a \emph{monotonically shrinking} Kleene
iteration starting above the true answer, \textbf{any operator we
substitute for an exact one, provided it only ever removes points that a
correct computation would also remove (i.e.\ computes a subset of the exact
result at every step), yields a sound under-approximation of the true
fixed point}: whatever survives to the end is genuinely in the true winning
region, so the final verdict is trustworthy whenever it is definitive; the
only possible defect is spurious non-convergence (falling back to
\textsc{Unknown}, or an unnecessary $K$-increment) where an exact
computation would have concluded. This one fact underwrites every
soundness argument in the remainder of the document.

\section{Detecting and Verifying the Symmetry Group}
\label{sec:detection}

We want a permutation group $\Phigroup$ acting on the states $\Q$ of $\Aut$
such that the $K$-bounded game is $\Phigroup$-equivariant, so that we may
canonicalize downsets under $\Phigroup$ without changing the winner. We
restrict attention to $\Phigroup$ generated by \emph{client-index
transpositions}: pairs of indices $a, b$ (read off atomic-proposition names
that share a prefix and differ only in a trailing integer, e.g.\ $r_a,
r_b$, $g_a, g_b, \dots$) such that swapping every occurrence of index $a$
and index $b$, across every such indexed family simultaneously, is a
structural automorphism of $\Aut$.

\begin{definition}[Structural automorphism]
Let $\pi$ be a permutation of the atomic propositions of $\Aut$ that maps
inputs to inputs and outputs to outputs, and let $\phi : \Q \to \Q$ be a
bijection. The pair $(\phi,\pi)$ is a \emph{structural automorphism} of
$\Aut$ if $\phi(q_0) = q_0$, and for every edge $q
\xrightarrow{c} q'$ of $\Aut$ there is an edge $\phi(q)
\xrightarrow{\pi(c)} \phi(q')$ with the same acceptance marking, and vice
versa.
\end{definition}

\begin{theorem}[Soundness of canonicalization under a verified automorphism]
\label{thm:auto-sound}
If $(\phi,\pi)$ is a structural automorphism of $\Aut$, then $\phi$ acts on
vectors coordinatewise ($(\phi \cdot v)(q) = v(\phi^{-1}(q))$), $\PreHat$ is
$\phi$-equivariant --- $\PreHat(\phi\cdot v,\, \pi(i),\, \pi(o)) = \phi \cdot
\PreHat(v, i, o)$ --- and consequently $\CPre$ is $\Phigroup$-equivariant
for $\Phigroup = \langle \phi \rangle$: $\CPre(\Phigroup \cdot F) =
\Phigroup \cdot \CPre(F)$ for every downset $F$. Since the initial safe
downset is $\Phigroup$-invariant and $\iota$ is fixed by every $\phi \in
\Phigroup$ (as $\phi(q_0)=q_0$), every iterate of the Kleene sequence
computing the greatest fixed point of $\CPre$, and hence the fixed point
itself, is $\Phigroup$-invariant.
\end{theorem}

The proof is immediate from the definitions and the fact that $\delta$ is
carried faithfully by $(\phi,\pi)$; it is exactly the standard automorphism
argument for symmetry reduction in model checking, specialized to this
particular backward operator. The practical content of
Theorem~\ref{thm:auto-sound} is that we never need to \emph{trust} a
candidate symmetry derived from, say, the surface syntax of the LTL
formula: we can and do \emph{verify it structurally on the automaton}
before using it, and discard any candidate that fails verification. Since
discarding a generator only shrinks $\Phigroup$, this is conservative by
construction (Remark~\ref{rem:conservative}).

\begin{remark}[Conservativeness of partial detection]
\label{rem:conservative}
If detection finds only a strict subgroup of the true (semantic) symmetry
group of the specification --- because, e.g., our syntactic candidate
generation misses a symmetry not reflected in AP naming, or because a
genuine symmetry fails our \emph{structural} (automaton-level) verification
even though it holds \emph{semantically} (language-level) --- the algorithm
still loses nothing in soundness, only in the amount of reduction achieved.
We never assume a symmetry we have not verified.
\end{remark}

\subsection{Candidate generation and structural verification}

Candidates are read off directly from AP names: every atomic proposition
matching \texttt{prefix\_index} (a maximal trailing run of digits) is
grouped by prefix into an indexed \emph{family}; the client-index set is
the intersection of the index sets of every family. For every pair of
client indices $(a,b)$ we test whether the transposition $\pi_{ab}$ ---
swap $a \leftrightarrow b$ simultaneously in every family --- is realized
by a structural automorphism, and if so, record the induced state
permutation $\phi_{ab}$ as a verified generator.

Because the game automaton $\Aut$ is in general \emph{non-deterministic}
(this is the case for every automaton arising in this pipeline, including
the small arbiter instances), we cannot simply propagate $\phi$ forward
from the initial state by following unique successor edges as one could
for a deterministic automaton. Instead we build $\mathcal B = \pi_{ab}(\Aut)$
by literally relabelling the atomic propositions of a copy of $\Aut$ (a
BDD-variable-pair substitution on every edge condition), and then search
for a \emph{label-preserving isomorphism} $\phi : \Aut \to \mathcal B$ ---
a bijection on states respecting edge labels, acceptance marks, and the
initial state --- via a backtracking search with signature-based pruning
(each state's signature is its acceptance bit, in-degree, and the sorted
multiset of (label,~acceptance) pairs on its outgoing edges; this is not
merely a heuristic filter but a necessary condition satisfied by any valid
$\phi$, so pruning on it never discards a correct assignment). If such a
$\phi$ exists, it is (by construction of $\mathcal B$) exactly a witness
that $(\phi,\pi_{ab})$ is a structural automorphism, and is recorded as a
generator; if the search fails, $\pi_{ab}$ is silently discarded.

We say $\Phigroup$ is \emph{full-symmetric} on the client-index set
$\{i_1,\dots,i_n\}$ if \emph{every} pairwise transposition $\pi_{ab}$
verifies. The Young-subgroup argument used later in
\S\ref{sec:cpre-step} depends on $\Phigroup$ being (at least) the full
symmetric group $S_n$ on client indices, not merely some subgroup that
happens to generate $S_n$ abstractly; we therefore gate all downstream use
on this stronger, directly-checked condition rather than on generating-set
closure.

\paragraph{Empirical confirmation on real automata.} On the arbiter family
($n=3,4,5,6$), detection recovers exactly the full symmetric group: $3/3$,
$6/6$, $10/10$, $15/15$ verified transpositions respectively, on the actual
non-deterministic game automata (not a simplified or synthetic model).

\section{Recovering the Client-Block Structure}
\label{sec:blocks}

Given a verified full-symmetric $\Phigroup = S_n$ on client indices, we
next need to know, for the purpose of representing vectors canonically,
\emph{which automaton states are ``owned'' by which client, and in which
structural role} --- i.e., the partition of $\Q$ into $B$ \emph{blocks} of
size $n$ (one state per client per role) plus a set of \emph{shared} states
fixed by every generator.

\begin{definition}[Orbit partition]
Let $\Gamma = \langle \phi_{ab} : (a,b)\ \text{verified} \rangle$ be the
permutation group on $\Q$ generated by the verified generators. The orbits
of $\Gamma$ acting on $\Q$ partition the automaton states; we require every
orbit to have size $1$ (a \emph{shared} state) or size $n$ (a
\emph{client-block} state), and reject (fall back to no symmetry
exploitation) if any other orbit size occurs.
\end{definition}

Orbits are computed by a standard union--find over the generator action
(apply every generator to every state and union the endpoints), which is
polynomial in $|\Q|$ and the number of generators ($O(n^2)$).

\subsection{Slot identity across blocks}

Knowing the orbit partition is not yet enough: to build a client's
\emph{type} (its value in every block simultaneously) we need a
\emph{consistent} labelling of ``which state in block $j$'' corresponds to
``the same client'' as ``which state in block $j'$''. Since automaton
states carry no inherent client label beyond what the group action
reveals, we recover this labelling via each state's \emph{stabilizer
signature}.

\begin{proposition}[Stabilizer signatures separate clients for $n\geq 3$]
\label{prop:stabilizer}
For $\Phigroup \cong S_n$ acting naturally on $n$ labelled points with
$n\geq 3$, the point-stabilizer of point $k$ --- the subgroup fixing $k$,
here identified with the \emph{subset of generators $\phi_{ab}$ that fix
$k$}, i.e.\ those with $k \notin \{a,b\}$ --- is a distinct subgroup for
each $k \in \{0,\dots,n-1\}$, and hence a distinct \emph{signature} (as a
bit-vector over the generator list). This fails for $n=2$: the unique
generator $\phi_{01}$ moves both points, so both have the identical (empty)
signature.
\end{proposition}

We therefore compute, for every state $s$ in a client-block, its signature
$\mathrm{sig}(s) = \{\, k : \phi_k(s) = s \,\}$ (indices into the verified
generator list), fix a canonical labelling on one reference block by
sorting its $n$ states by raw index, and match every other block's states
to the reference by exact signature equality. We \emph{verify at run time}
that (a) the reference block's $n$ signatures are pairwise distinct, (b)
every other block has exactly one state per reference signature, and
decline (return ``no usable layout'') if either check fails --- in
particular, and by design, whenever $n=2$ specifically. This is again a
conservative choice: $n=2$ is not the regime this optimization targets
(Table in \S\ref{sec:orbit-positive} shows the blow-up only bites for
$n\gtrsim 5$), so we decline rather than engineer a bespoke, harder-to-verify
tie-break for a case of no practical interest here.

\paragraph{Validation.} The algorithm was validated two ways. First,
synthetically: on hand-built groups emulating $(n,B,S)$-structured
automata for $n\in\{2,3,5,7\}$, $B\in\{1,3,4\}$, $S\in\{0,2\}$ (24
configurations), the recovered layout matches the intended structure
exactly in every $n\geq3$ case, and $n=2$ correctly declines. Second,
against the \emph{real} arbiter automata's detected generators (not
synthetic): the recovered layout is $4$ blocks, $3$ shared states for
every $n=3,4,5,6$ instance --- matching, exactly, an entirely independent
offline analysis of the antichain's raw column signatures
(\S\ref{sec:orbit-positive}). Two unrelated analyses --- one on automaton
structure, one on antichain data --- agreeing exactly is strong
corroborating evidence that the recovered block structure is the correct
one.

\section{Canonical Representation: Count-Vectors}
\label{sec:count-vectors}

\begin{definition}[Type, count-vector]
Fix a client-block layout with $B$ blocks and shared coordinates $S$. A
client's \emph{type} is its tuple of values across the $B$ blocks, an
element of $D = \{-1,\dots,K\}^B$; write $r = |D| \leq (K+2)^B$. A
\emph{count-vector} is a pair $(s, c)$ where $s \in \{-1,\dots,K\}^{|S|}$ is
the shared part and $c : D \to \mathbb N$ records, for each type, how many
of the $n$ clients realize it ($\sum_{t} c(t) = n$).
\end{definition}

\begin{proposition}[Count-vectors are exactly $\Phigroup$-orbit representatives]
Two raw vectors $u, v$ (agreeing on the shared part) lie in the same
$\Phigroup$-orbit if and only if they induce the same count-vector.
\end{proposition}

This is immediate: $\Phigroup = S_n$ acts by arbitrarily permuting which
client owns which type, so the orbit of a vector is exactly determined by
the \emph{multiset} of client types, i.e.\ by $c$. Crucially, $r$ is
bounded by a constant \emph{independent of $n$} (it depends only on $B$
and the current $K$), so a count-vector has boundedly many nonzero entries
regardless of how many clients there are --- this is the representation
that achieves the linear orbit count of \S\ref{sec:orbit-positive}, in
contrast to the raw antichain (exponential in $n$) and the BDD encoding
(also exponential in $n$, \S\ref{sec:bdd-negative}).

\section{Domination Algebra on Count-Vectors}
\label{sec:algebra}

Acacia's downset representation needs three operations: membership
(\emph{contains}), \emph{union}, and \emph{intersect}. We derive each for
count-vectors.

\subsection{Exact: \textsc{Contains} via bipartite transportation feasibility}

\begin{proposition}[Domination as a transportation problem]
\label{prop:contains}
A count-vector $u=(s_u,c_u)$ dominates $v=(s_v,c_v)$ (i.e.\ some raw
realization of $u$ dominates some/every raw realization of $v$ under the
same client-index assignment) iff $s_u \dom s_v$ pointwise \emph{and} there
is a feasible transportation plan from supply $c_u$ to demand $c_v$ using
only edges $(t,t')$ with $t \dom t'$ (pointwise domination of type-tuples).
\end{proposition}

\begin{proof}[Proof sketch]
A flow of value $n$ through the bipartite network with left nodes $D$
(capacity $c_u(t)$), right nodes $D$ (capacity $c_v(t')$), and edges
exactly $\{(t,t') : t \dom t'\}$ decomposes into unit flows, each
specifying that one client realizing $u$-type $t$ is paired with one
client realizing $v$-type $t'$ with $t \dom t'$; conservation at every node
guarantees every $u$-client and every $v$-client is matched exactly once.
Such a decomposition exists iff the network admits a flow saturating all
demand, i.e.\ iff $\max\text{-flow} = n$ (since $\sum c_u = \sum c_v = n$).
Conversely, a max flow of value $n$ gives exactly such a pairing, i.e.\ a
bijection between the $n$ clients of one realization and the $n$ clients of
the other, respecting the required pointwise domination client-by-client
--- precisely the definition of raw vector domination.
\end{proof}

This is exact and polynomial in $r$ (independent of $n$): the network has
$O(r)$ nodes, so a standard max-flow algorithm (we use a simple
Edmonds--Karp-style augmenting-path search, pushing full bottleneck
capacity per augmentation, giving $O(r)$ augmentations independent of the
magnitude of $n$) suffices. \textsc{Union} of two count-vector antichains
reduces to pairwise dominance filtering with the same test --- no
additional combinatorics is required, since the antichain-of-maximal-
elements of a union of downsets is simply the union of the two antichains
with dominated elements removed.

\begin{remark}[Validation]
\textsc{Contains} and \textsc{Union} were validated by $1500$ randomized
trials cross-checked against brute-force enumeration of all $n!$
permutation pairings for small $n$ (up to $6$--$7$ clients, $B$ up to $4$,
$K$ up to $4$): $100\%$ agreement.
\end{remark}

\subsection{Inexact by necessity: \textsc{Intersect}}
\label{sec:intersect}

The classical antichain identity
\[
  \down{A} \cap \down{B} = \down{\{\, \mathrm{meet}(a,b) : a \in A,\ b \in B \,\}}
\]
still holds for count-vectors, but the set of achievable meet outcomes
from pairing two type-distributions \emph{depends on which
client-to-client pairing (transportation plan) is used} --- unlike
\textsc{Contains}, which only asks whether \emph{some} pairing exists,
\textsc{Intersect} needs the \emph{maximal achievable} meet outcomes over
\emph{all} pairings, i.e.\ effectively the extreme points of the
transportation polytope between $c_u$ and $c_v$. We were unable to find a
polynomial (in $r$) characterization of this set; the number of extreme
points of a transportation polytope is not bounded by a polynomial in the
number of supply/demand nodes in general.

\subsubsection{The resolution: soundness does not require completeness}

\begin{lemma}[Under-approximating \textsc{Intersect} is sound]
\label{lem:intersect-sound}
Let $I(A,B)$ be any procedure that returns a subset of the exact
antichain-of-maximal-elements of $\down A \cap \down B$, where every
returned point is a genuinely achievable meet (i.e.\ realized by some
valid transportation plan between some $a\in A$ and some $b\in B$). Then
substituting $I$ for exact intersection throughout Acacia's Kleene
iteration (\S\ref{sec:background}) preserves soundness: every
\textsc{Realizable}/\textsc{Unrealizable} verdict reached remains correct.
\end{lemma}

\begin{proof}
By the greatest-fixed-point monotonicity argument of
\S\ref{sec:background}: since $I(A,B) \subseteq \down A \cap \down B$ at
every step, the computed sequence of iterates is pointwise $\subseteq$ the
exact Kleene sequence, hence the computed fixed point is a subset of the
true winning region. Anything that survives (in particular the initial
vector $\iota$) is therefore genuinely winning. The only possible defect
is the computed fixed point excluding points the exact computation would
have kept, which can only cause spurious non-termination or an
unnecessary $K$-increment, never a false positive.
\end{proof}

\subsubsection{A concrete, bounded construction}

We generate candidate merges via the \emph{Northwest-corner} construction
from classical transportation theory: sort both type-multisets by a chosen
key and greedily match largest-with-largest (or by any other fixed
orientation), producing, by construction, \emph{one} valid transportation
plan of size $O(r)$ in $O(r\log r)$ time. We use four orientations (sum
ascending/descending on each side) per antichain-element pair, giving a
bounded ($O(1)$ per pair) set of candidates; each is, individually, a
provably valid achievable point (Lemma~\ref{lem:intersect-sound}'s
hypothesis is satisfied trivially, by construction, for every candidate we
emit), so soundness is structural and requires no separate proof per
candidate --- only \emph{which} maximal points are found is heuristic.

\begin{remark}[Validation]
Soundness (never emitting a point outside the true intersection) was
validated by $400$ single-pair and $80$ multi-element-antichain randomized
trials, cross-checked against brute-force enumeration of the true
intersection antichain (via all $n!$ pairings): \emph{zero} unsound points
across all trials. Completeness (recovering literally every true maximal
point) is, as expected, partial: roughly $45$--$60\%$ of true maximal
points are recovered exactly, which is immaterial to correctness by
Lemma~\ref{lem:intersect-sound} but is recorded as a quality signal.
\end{remark}

\section{Assembling the Symmetric \CPre\ Step}
\label{sec:cpre-step}

We now assemble \S\ref{sec:algebra}'s operations into a replacement for
one iteration of $\CPre$ that keeps $F$ in canonical (count-vector,
$\Phigroup$-invariant) form throughout, never materializing a raw
antichain or enumerating $\Phigroup$.

\subsection{Representative inputs via Young subgroups}

For a specific raw input $i$ (say, ``exactly $k$ of $n$ clients request''),
its stabilizer $\Stab_\Phigroup(i) \leq \Phigroup$ is the \emph{Young
subgroup} $S_k \times S_{n-k}$: permutations may freely rearrange the $k$
requesting clients among themselves, and freely rearrange the $n-k$
non-requesting clients among themselves, but may not move a client across
the two groups. Consequently the orbit of $i$ under $\Phigroup$ is exactly
``every input with the same requesting count $k$'', and there are only
$n+1$ such orbits ($k=0,\dots,n$) --- linear in $n$, as opposed to the
$2^n$ raw inputs.

\subsection{The obstacle: no canonical split of a canonical vector's clients}

Fix a representative input $i$ for orbit $k$. We want
\[
  T_i \;=\; \bigcup_{o\ \mathrm{compat.}\ i} \PreHat(F, i, o),
\]
computed from $F$'s canonical (unsplit) antichain. The obstacle is that
$\PreHat(\cdot,i,o)$, for the \emph{fixed}, non-symmetric $i$, is only
$\Stab_\Phigroup(i)$-equivariant, not $\Phigroup$-equivariant (a
straightforward adaptation of the proof of Theorem~\ref{thm:auto-sound},
restricting to $\phi \in \Stab_\Phigroup(i)$). A canonical $u \in F$ does
not itself record \emph{which} of its (interchangeable, same-typed)
clients would be ``the $k$ requesting ones'' were we to realize a concrete
raw vector and apply $i$ to it --- different choices of which clients play
the requesting role can, in general, yield genuinely different $\PreHat$
outcomes, so there is no single canonical answer.

\subsection{The resolution: reuse the \textsc{Intersect} soundness argument}

Since $T_i$ only ever appears as a factor intersected into $F$, exactly
the same reasoning as Lemma~\ref{lem:intersect-sound} applies: an
under-approximated $T_i$ (any subset of the true $T_i$ built entirely from
genuinely achievable points) keeps the overall computation sound. We
therefore build $T_i$ from a \emph{bounded set of candidate
client-to-role assignments} per canonical $u$ --- concretely, sort $u$'s
clients by a small fixed family of keys (total value, sum of the first
coordinate, ascending and descending) and assign the top $k$ to the
requesting role --- realize each candidate as one concrete raw vector,
apply the \emph{existing, unmodified} backward action $\PreHat$
(\texttt{actioners::standard::apply}, reused verbatim: no new game
semantics are introduced or need re-verification), convert the raw result
back to canonical form, and take the union (\S\ref{sec:algebra}) over all
candidates and all compatible outputs $o$.

\begin{proposition}[Un-splitting is exact, not approximate]
Converting a raw result vector back to an (unsplit) canonical count-vector
via the type-counting map of \S\ref{sec:count-vectors} loses no
information relative to any intermediate ``which clients were assigned the
requesting role'' bookkeeping: that assignment is external scaffolding
used only to construct one valid raw vector, and is never part of the raw
vector's coordinates.
\end{proposition}

\subsection{The resulting algorithm}

\begin{algorithm}[H]
\caption{One symmetric \CPre\ step for representative input $i$ (orbit
size $k$), replacing \texttt{cpre\_inplace}}
\label{alg:cpre-step}
\begin{algorithmic}[1]
\Require canonical antichain $F$ (list of count-vectors); block layout $L$;
   representative input $i$; its compatible outputs $O_i$
\Ensure a sound under-approximation of $F \cap \bigcup_{o\in O_i}\PreHat(F,i,o)$,
   canonical
\State $C \gets \emptyset$ \Comment{candidate merge points for $T_i$}
\ForAll{$u \in F$}
  \ForAll{split key $\kappa \in \{\text{sum}\!\uparrow, \text{sum}\!\downarrow,
           \text{1st-coord}\!\uparrow, \text{1st-coord}\!\downarrow\}$}
    \State $v_{\mathrm{raw}} \gets \textsc{Realize}(u, L, \kappa)$
       \Comment{one concrete raw vector for $u$, split by $\kappa$}
    \ForAll{$o \in O_i$}
      \State $r_{\mathrm{raw}} \gets \textsc{ActionerApply}(v_{\mathrm{raw}}, i, o, \mathrm{backward})$
         \Comment{unmodified existing operator}
      \State $C \gets C \cup \{\textsc{ToCount}(r_{\mathrm{raw}}, L)\}$
    \EndFor
  \EndFor
\EndFor
\State $T_i \gets \textsc{Union}(C)$ \Comment{exact, \S\ref{sec:algebra}}
\State \Return $\textsc{Intersect}(F, T_i)$ \Comment{capped-sound, \S\ref{sec:intersect}}
\end{algorithmic}
\end{algorithm}

The outer solve loop (mirroring \texttt{k\_bounded\_safety\_aut::solve})
iterates Algorithm~\ref{alg:cpre-step} over the $n+1$ representative
inputs until none changes $F$, checks $\textsc{Contains}(F,\iota)$
(\S\ref{sec:algebra}, exact) for convergence, and increments $K$ exactly as
the existing solver does (shifting the counting-coordinate entries of
every type in every count-vector by the increment; no new algebra is
needed for this step).

\begin{theorem}[End-to-end soundness]
The algorithm obtained by replacing \texttt{cpre\_inplace} with
Algorithm~\ref{alg:cpre-step}, run to the same convergence criterion as
the existing solver, never returns \textsc{Realizable} for an
unrealizable specification nor \textsc{Unrealizable} for a realizable one.
It may, relative to the exact antichain algorithm, answer
\textsc{Unknown} (or perform an unnecessary $K$-increment) on instances
the exact algorithm would have resolved.
\end{theorem}

\begin{proof}
Compose Theorem~\ref{thm:auto-sound} (canonicalization under a verified
automorphism is exact and loses nothing) with
Lemma~\ref{lem:intersect-sound} applied twice: once directly to
\textsc{Intersect} in the final line of Algorithm~\ref{alg:cpre-step}, and
once to the construction of $T_i$ itself (a union of genuinely achievable
points is, a fortiori, a subset of the true $T_i$, and this subset relation
is exactly what Lemma~\ref{lem:intersect-sound}'s hypothesis needs). Both
substitutions are one-directional (subset-of-true), so their composition
is as well; the greatest-fixed-point argument of \S\ref{sec:background}
then gives the result.
\end{proof}

\section{Implementation Status and Remaining Work}
\label{sec:status}

The branch now contains both the core symmetry modules and an opt-in live
decision path. The main components are all under \texttt{src/solver/}, with
corresponding tests under the \texttt{symmetry} \texttt{meson} test suite:

\begin{itemize}[nosep]
  \item \texttt{symmetry.hh} --- detection and structural verification of
    client-transposition automorphisms (\S\ref{sec:detection}); confirmed
    to recover the full symmetric group on real, non-deterministic arbiter
    automata for $n=3,\dots,6$.
  \item \texttt{symmetric\_blocks.hh} --- client-block and slot recovery
    via stabilizer signatures (\S\ref{sec:blocks}); validated on $24$
    synthetic configurations and cross-validated against the real arbiter
    automata's detected generators, agreeing exactly with an independent
    offline analysis.
  \item \texttt{symmetric\_downset.hh} --- the count-vector domination
    algebra (\S\ref{sec:algebra}): exact \textsc{Contains} and
    \textsc{Union}, capped-sound \textsc{Intersect}; validated by
    thousands of randomized cross-checks against brute-force ground
    truth, with zero soundness violations observed.
  \item \texttt{symmetric\_conversion.hh} --- raw-vector $\leftrightarrow$
    count-vector conversion and the bounded candidate-split key family
    used in Algorithm~\ref{alg:cpre-step}; validated by exhaustive
    round-trip tests.
  \item \texttt{symmetric\_k\_bounded\_safety\_aut.hh} --- the live
    quotient fixed-point loop, representative input/output dispatch, bounded
    \texttt{union\_o} construction, and fallback boundary to the existing
    solver. This path is decision-only and guarded by
    \texttt{ACACIA\_ENABLE\_SYMMETRIC\_SOLVER}.
  \item \texttt{symmetric\_dense\_downset.hh} and
    \texttt{symmetric\_profile.hh} --- the current dense/SIMD-oriented
    adapter and profiling hooks used to measure and optimize the quotient
    hot path. Dense union is exact-equivalent to the sparse implementation;
    dense intersection preserves the same bounded-sound semantics as the
    sparse path.
\end{itemize}

\noindent The live solver integration is now present but still
experimental. It is default-off, bounded by explicit work caps, and falls
back to the existing antichain solver whenever symmetry is absent,
representative construction is capped, or the quotient loop does not reach a
trusted decision cheaply enough. Targeted TLSF arbiter benchmarking against
\texttt{origin/master} showed matching verdict coverage and no polarity
mismatches, but no performance win yet; the current bottleneck is still the
quotient predecessor/intersection hot path around larger $K$ values.
Strategy extraction for controller synthesis (as opposed to the
realizability decision) still requires an additional
canonical-representative tie-break to reconstruct a concrete,
symmetry-\emph{breaking} strategy from the quotient computation, and remains
out of scope for this document.

\end{document}
